
Semantic Role Labelling~\citep[SRL, ][]{marquez08srl} is generally understood as 
the task of identifying and classifying the semantic arguments and modifiers of 
the predicates mentioned in a sentence. For example, in the case of the 
following sentence\footnote{``Haag plays Elianti'' is a segment of a sentence in 
training corpus.}:
\begin{quote}
\begin{center}
    \includegraphics[scale=.63]{haag-example}
\end{center}
\end{quote}
we are to find out that for the predicate token {}``plays'' with sense ``play a 
role'' (play.02) the phrase headed by the token {}``Haag'' is referring to the 
player (A0) of the play event, and the phrase headed by the token {}``Elianti''  
is referring to the role (A1) being played. SRL is considered as a key task for 
applications that require to answer {}``Who'', {}``What'', {}``Where'', etc.  
questions, such as Information Extraction, Question Answering and Summarization.  

In this paper we introduce a Markov Logic~\citep[ML,][]{richardson06mln} approach to multi-lingual
SRL. We present a brief introduction to ML in 
section \ref{sec:markovlogic}.  At the core of ML are Markov Logic Networks~(MLN): sets of weighted First Order 
Logic (FOL) formulae. One attractive feature of the ML 
framework is that we can divide FOL formulae into subsets that serve as re-usable ``modules''. 
With this in mind we define MLN modules which capture the different stages of 
the task: argument identification, argument classification, and sense 
disambiguation. We also define modules for some of the language specific aspects of 
the task. We present all our ML modules in section 
\ref{sec:model}. These modules are the building blocks that we use to create 
MLNs for specific languages. We present our experiments and results for each of 
the languages of the task in section \ref{sec:results}. In section 
\ref{sec:analysis} we present some error analysis.


